% ============================================================
% Chapter 8: Identity, Directory Services and Federated Single Sign-On
% Language: Greek — edit freely; regenerate from src/content if preferred.
% ============================================================
\chapter{Ταυτότητα, Υπηρεσίες Καταλόγου και Ομοσπονδιακή Ενιαία Σύνδεση}\label{ch:8}
\chaptermeta{IAM · Παράγοντες ταυτοποίησης και MFA · FIDO2 · RBAC και ABAC · LDAP και Active Directory · Kerberos · SAML, OAuth 2.0 και OpenID Connect}{Επίπεδο: Μεσαίο επίπεδο · Ταυτότητα}{Χρόνος μελέτης: 9–11 ώρες}

Καθώς οι οργανισμοί μεταφέρονται σε υπηρεσίες νέφους και στην απομακρυσμένη εργασία, η περίμετρος του δικτύου χάνει το νόημά της και \textbf{η ταυτότητα γίνεται η νέα περίμετρος}. Οι περισσότερες σύγχρονες εισβολές δεν σπάνε την κρυπτογράφηση ούτε εκμεταλλεύονται εξωτικές ευπάθειες· συνδέονται με κλεμμένα ή καταχρηστικά χρησιμοποιούμενα διαπιστευτήρια. Η κατανόηση του τρόπου με τον οποίο δημιουργούνται, επαληθεύονται, εξουσιοδοτούνται και ομοσπονδοποιούνται οι ταυτότητες είναι επομένως κεντρική για τη σύγχρονη άμυνα.

Το κεφάλαιο στηρίζεται στα κρυπτογραφικά θεμέλια του Κεφαλαίου 7. Ξεκινά με την αρχιτεκτονική της διαχείρισης ταυτότητας και πρόσβασης και με τους παράγοντες που χρησιμοποιούνται για την ταυτοποίηση, συμπεριλαμβανομένου του ανθεκτικού στο ψάρεμα FIDO2. Στη συνέχεια συγκρίνει τα μοντέλα ελέγχου πρόσβασης, εξηγεί τις υπηρεσίες καταλόγου και το πρωτόκολλο Kerberos στο οποίο στηρίζονται οι τομείς Windows —μαζί με τις επιθέσεις εναντίον του— και ολοκληρώνεται με τα πρωτόκολλα ομοσπονδίας SAML, OAuth 2.0 και OpenID Connect, που επιτρέπουν την ενιαία σύνδεση πέρα από τα όρια ενός οργανισμού.

\begin{outcomesbox}{Μαθησιακά αποτελέσματα}
Μετά τη μελέτη του κεφαλαίου θα πρέπει να μπορείτε:
\begin{itemize}
  \item Να διακρίνετε την αναγνώριση, την ταυτοποίηση (αυθεντικοποίηση), την εξουσιοδότηση και την καταγραφή.
  \item Να συγκρίνετε τους παράγοντες ταυτοποίησης και να εξηγείτε γιατί το FIDO2 αντιστέκεται στο ηλεκτρονικό ψάρεμα.
  \item Να επιλέγετε μεταξύ RBAC, ABAC και ελέγχου πρόσβασης βάσει κανόνων για ένα δεδομένο σενάριο.
  \item Να εξηγείτε τη ροή εισιτηρίων του Kerberos και τις αρχές πίσω από τις επιθέσεις Kerberoasting, pass-the-ticket και golden ticket.
  \item Να περιγράφετε τον ρόλο των SAML, OAuth 2.0 και OpenID Connect, καθώς και συνήθεις αδυναμίες της ομοσπονδίας ταυτοτήτων.
\end{itemize}
\end{outcomesbox}

\section{Αρχιτεκτονική IAM: αναγνώριση, ταυτοποίηση, εξουσιοδότηση και καταγραφή}\label{sec:8.1}
Η \textbf{διαχείριση ταυτότητας και πρόσβασης} (Identity and Access Management, IAM) οργανώνεται γύρω από τέσσερα διακριτά βήματα. Η \textbf{αναγνώριση} είναι η δήλωση μιας ταυτότητας, όπως η πληκτρολόγηση ενός ονόματος χρήστη. Η \textbf{ταυτοποίηση} ή \textbf{αυθεντικοποίηση} είναι η επαλήθευση αυτής της δήλωσης, για παράδειγμα με έλεγχο ενός κωδικού ή ενός κλειδιού ασφαλείας. Η \textbf{εξουσιοδότηση} καθορίζει τι επιτρέπεται να κάνει το υποκείμενο που ταυτοποιήθηκε. Η \textbf{καταγραφή} (ή λογιστική παρακολούθηση) καταγράφει τι έκανε πράγματι το υποκείμενο, υποστηρίζοντας τη λογοδοσία και τις έρευνες. Η σύγχυση αυτών των βημάτων οδηγεί σε σχεδιαστικά σφάλματα —για παράδειγμα, στην αντιμετώπιση ενός επιτυχώς ταυτοποιημένου χρήστη ως αυτομάτως εξουσιοδοτημένου για κάθε πόρο.

Η IAM καλύπτει επίσης τον \textbf{κύκλο ζωής της ταυτότητας}: οι λογαριασμοί πρέπει να δημιουργούνται όταν κάποιος εντάσσεται στον οργανισμό, να προσαρμόζονται όταν αλλάζει ρόλο και να απενεργοποιούνται άμεσα όταν αποχωρεί. Οι ορφανοί λογαριασμοί πρώην εργαζομένων και η σταδιακή συσσώρευση δικαιωμάτων κατά τη διάρκεια μιας επαγγελματικής πορείας —η λεγόμενη \emph{διολίσθηση προνομίων} (privilege creep)— συγκαταλέγονται στα πιο συνηθισμένα ευρήματα των ελέγχων ασφάλειας. Το καθιερωμένο αντίμετρο είναι οι περιοδικές \textbf{επισκοπήσεις πρόσβασης}, κατά τις οποίες οι προϊστάμενοι επιβεβαιώνουν ότι κάθε δικαίωμα εξακολουθεί να είναι αναγκαίο.

\section{Παράγοντες ταυτοποίησης, MFA και FIDO2 ανθεκτικό στο ψάρεμα}\label{sec:8.2}
Οι παράγοντες ταυτοποίησης ομαδοποιούνται παραδοσιακά σε τρεις κατηγορίες: \textbf{κάτι που γνωρίζεις} (κωδικός ή PIN), \textbf{κάτι που έχεις} (κινητό τηλέφωνο, έξυπνη κάρτα ή κλειδί ασφαλείας) και \textbf{κάτι που είσαι} (βιομετρικό χαρακτηριστικό, όπως το δακτυλικό αποτύπωμα). Ο \textbf{έλεγχος ταυτότητας πολλαπλών παραγόντων} (Multi-Factor Authentication, MFA) συνδυάζει παράγοντες από \emph{διαφορετικές} κατηγορίες, ώστε η υποκλοπή ενός παράγοντα να μην αρκεί. Ωστόσο, δεν είναι όλες οι μορφές MFA εξίσου ισχυρές. Οι κωδικοί μίας χρήσης μέσω SMS μπορούν να υποκλαπούν με αντικατάσταση κάρτας SIM (SIM swapping)· οι χρονικά μεταβαλλόμενοι κωδικοί εφαρμογών ταυτοποίησης (TOTP, σύμφωνα με τα πρότυπα OATH) και οι ειδοποιήσεις push μπορούν ακόμη να αναμεταδοθούν από έναν διακομιστή ψαρέματος σε πραγματικό χρόνο.

Το \textbf{FIDO2/WebAuthn} αλλάζει ριζικά αυτή την εικόνα. Κατά την εγγραφή, ο αυθεντικοποιητής —ένα κλειδί ασφαλείας ή το ασφαλές υλικό ενός κινητού ή φορητού υπολογιστή— δημιουργεί ένα ζεύγος κλειδιών ειδικά για τον συγκεκριμένο ιστότοπο. Κατά τη σύνδεση, υπογράφει μια πρόκληση που περιλαμβάνει την προέλευση (origin) του ιστοτόπου. Ένας ψεύτικος ιστότοπος με διαφορετικό όνομα τομέα λαμβάνει επομένως μια υπογραφή άχρηστη για τον πραγματικό ιστότοπο, και δεν υπάρχει κοινό μυστικό προς υποκλοπή. Γι' αυτό το FIDO2 χαρακτηρίζεται \textbf{ανθεκτικό στο ηλεκτρονικό ψάρεμα} και γι' αυτό τα \emph{passkeys} που βασίζονται σε αυτό αντικαθιστούν τους κωδικούς σε πολλές υπηρεσίες.

\section{Μοντέλα ελέγχου πρόσβασης: RBAC, ABAC και έλεγχος βάσει κανόνων}\label{sec:8.3}
Αφού ταυτοποιηθεί ένα υποκείμενο, το σύστημα πρέπει να αποφασίσει τι επιτρέπεται να κάνει. Στον \textbf{έλεγχο πρόσβασης βάσει ρόλων} (Role-Based Access Control, RBAC), τα δικαιώματα ανατίθενται σε ρόλους (για παράδειγμα \emph{νοσηλευτής}, \emph{λογιστής}, \emph{διαχειριστής βάσης δεδομένων}) και οι χρήστες λαμβάνουν ρόλους. Το RBAC είναι εύκολο στην κατανόηση και στον έλεγχο, όμως οι μεγάλοι οργανισμοί μπορεί να αντιμετωπίσουν \emph{έκρηξη ρόλων} όταν κάθε εξαίρεση απαιτεί νέο ρόλο. Ο \textbf{έλεγχος πρόσβασης βάσει χαρακτηριστικών} (Attribute-Based Access Control, ABAC) αξιολογεί πολιτικές επί χαρακτηριστικών του υποκειμένου, του πόρου, της ενέργειας και του περιβάλλοντος —για παράδειγμα, «οι ιατροί μπορούν να διαβάζουν τους φακέλους ασθενών της δικής τους κλινικής, κατά τη διάρκεια της βάρδιάς τους, από διαχειριζόμενη συσκευή». Το ABAC είναι πιο εκφραστικό, αλλά δυσκολότερο στην ανάλυση και στον έλεγχο.

Ο \textbf{έλεγχος πρόσβασης βάσει κανόνων} εφαρμόζει γενικούς κανόνες που δεν εξαρτώνται από τον μεμονωμένο χρήστη, όπως οι κανόνες ενός τείχους προστασίας ή ο κανόνας «καμία σύνδεση μεταξύ μεσονυκτίου και 5 π.μ.». Στην πράξη οι οργανισμοί συνδυάζουν τα μοντέλα: οι ρόλοι παρέχουν μια σταθερή βάση, τα χαρακτηριστικά εξειδικεύουν τις αποφάσεις ανάλογα με το πλαίσιο και οι γενικοί κανόνες επιβάλλουν περιορισμούς για ολόκληρο τον οργανισμό.

\section{Υπηρεσίες καταλόγου: X.500, LDAP και Active Directory}\label{sec:8.4}
Μια \textbf{υπηρεσία καταλόγου} είναι μια εξειδικευμένη, ιεραρχική βάση δεδομένων που αποθηκεύει πληροφορίες για χρήστες, ομάδες, υπολογιστές και άλλους πόρους και είναι βελτιστοποιημένη για συχνές αναγνώσεις. Το εννοιολογικό της μοντέλο προέρχεται από τα πρότυπα \textbf{X.500}, ενώ το \textbf{Lightweight Directory Access Protocol (LDAP)} είναι ο συνήθης τρόπος υποβολής ερωτημάτων και τροποποίησης. Κάθε εγγραφή προσδιορίζεται από ένα \emph{διακεκριμένο όνομα} (distinguished name), όπως \texttt{CN=Maria Papadopoulou,OU=Finance,DC=example,DC=com}.

Οι \textbf{Active Directory Domain Services (AD DS)} της Microsoft συνδυάζουν έναν κατάλογο LDAP με ταυτοποίηση Kerberos, DNS και πολιτικές ομάδας (Group Policy). Οργανώνουν τα αντικείμενα σε \emph{τομείς}, \emph{δέντρα} και \emph{δάση}, ενώ οι \emph{ελεγκτές τομέα} αναπαράγουν τον κατάλογο μεταξύ τους. Επειδή το AD ελέγχει την ταυτοποίηση σχεδόν κάθε συστήματος σε μια επιχείρηση Windows, η παραβίαση ενός λογαριασμού διαχειριστή τομέα σημαίνει συνήθως παραβίαση ολόκληρου του οργανισμού. Γι' αυτό η ασφάλεια του AD οργανώνεται γύρω από την \textbf{κλιμακωτή διαχείριση} (tiered administration): οι λογαριασμοί υψηλών προνομίων χρησιμοποιούνται μόνο σε αποκλειστικά, θωρακισμένα συστήματα και ποτέ σε συνηθισμένους σταθμούς εργασίας, όπου τα διαπιστευτήριά τους θα μπορούσαν να υποκλαπούν.

\section{Το Kerberos και οι επιθέσεις εναντίον του}\label{sec:8.5}
Το \textbf{Kerberos} επιτρέπει στους χρήστες να ταυτοποιούνται μία φορά και στη συνέχεια να έχουν πρόσβαση σε πολλές υπηρεσίες χωρίς να αποστέλλουν τον κωδικό τους μέσω του δικτύου. Ένα αξιόπιστο \textbf{Κέντρο Διανομής Κλειδιών} (Key Distribution Center, KDC) —στο AD, κάθε ελεγκτής τομέα— περιλαμβάνει δύο λογικές υπηρεσίες. Κατά τη σύνδεση, η \textbf{Υπηρεσία Ταυτοποίησης} (Authentication Service, AS) επαληθεύει τον χρήστη και εκδίδει ένα \textbf{Εισιτήριο Χορήγησης Εισιτηρίων} (Ticket-Granting Ticket, TGT), κρυπτογραφημένο με το κλειδί του ειδικού λογαριασμού \texttt{krbtgt}. Όταν ο χρήστης θέλει να προσπελάσει μια υπηρεσία, ο πελάτης παρουσιάζει το TGT στην \textbf{Υπηρεσία Χορήγησης Εισιτηρίων} (Ticket-Granting Service, TGS), η οποία εκδίδει ένα \emph{εισιτήριο υπηρεσίας} κρυπτογραφημένο με το κλειδί του λογαριασμού που εκτελεί την υπηρεσία. Η υπηρεσία αποκρυπτογραφεί το εισιτήριο και παραχωρεί την πρόσβαση. Τα εισιτήρια έχουν περιορισμένη διάρκεια ισχύος και το πρωτόκολλο προϋποθέτει συγχρονισμένα ρολόγια.

Ο σχεδιασμός αυτός επιτρέπει αρκετές γνωστές επιθέσεις. Στο \textbf{Kerberoasting}, οποιοσδήποτε χρήστης του τομέα ζητά εισιτήρια υπηρεσίας για λογαριασμούς που διαθέτουν \emph{όνομα κύριας υπηρεσίας} (Service Principal Name, SPN) και τα «σπάει» εκτός σύνδεσης· οι αδύναμοι κωδικοί λογαριασμών υπηρεσιών υποχωρούν γρήγορα. Στο \textbf{pass-the-ticket}, κλεμμένα εισιτήρια επαναχρησιμοποιούνται από άλλο μηχάνημα. Ένα \textbf{golden ticket} είναι ένα πλαστό TGT που δημιουργείται με το κλεμμένο κλειδί του \texttt{krbtgt} και παρέχει απεριόριστη πρόσβαση στον τομέα. Τα αντίμετρα περιλαμβάνουν μεγάλους τυχαίους κωδικούς ή διαχειριζόμενους λογαριασμούς υπηρεσιών ομάδας (gMSA), χρήση αποκλειστικά κρυπτογράφησης AES, προστασία των ελεγκτών τομέα ως στοιχείων της ανώτατης βαθμίδας (Tier 0), διπλή επαναφορά του κωδικού του \texttt{krbtgt} μετά από παραβίαση και παρακολούθηση ασυνήθιστων αιτημάτων εισιτηρίων.

\section{Ομοσπονδιακή ταυτότητα και ενιαία σύνδεση: SAML, OAuth 2.0 και OpenID Connect}\label{sec:8.6}
Η \textbf{ομοσπονδία ταυτοτήτων} (federation) επεκτείνει την ενιαία σύνδεση πέρα από έναν οργανισμό: ο χρήστης ταυτοποιείται στον δικό του \textbf{πάροχο ταυτότητας} (Identity Provider, IdP) και πολλοί \textbf{πάροχοι υπηρεσιών} αποδέχονται αυτή την ταυτοποίηση. Το \textbf{SAML 2.0} ανταλλάσσει ψηφιακά υπογεγραμμένους \emph{ισχυρισμούς} (assertions) σε XML και χρησιμοποιείται ευρέως σε επιχειρησιακές εφαρμογές ιστού. Το \textbf{OAuth 2.0} δεν είναι, αυστηρά μιλώντας, πρωτόκολλο ταυτοποίησης αλλά πλαίσιο \emph{εξουσιοδότησης}: επιτρέπει σε έναν χρήστη να παραχωρήσει σε μια εφαρμογή περιορισμένη πρόσβαση στους πόρους του (για παράδειγμα «ανάγνωση του ημερολογίου μου») μέσω έκδοσης ενός \emph{διακριτικού πρόσβασης} (access token), χωρίς να κοινοποιήσει τον κωδικό του. Το \textbf{OpenID Connect (OIDC)} προσθέτει ένα επίπεδο ταυτότητας πάνω από το OAuth 2.0, μέσω ενός υπογεγραμμένου \emph{διακριτικού ταυτότητας} (ID token) που δηλώνει ποιος είναι ο χρήστης.

Η ομοσπονδία συγκεντρώνει την εμπιστοσύνη και, επομένως, τον κίνδυνο. Αν ένας επιτιθέμενος υποκλέψει το κλειδί υπογραφής διακριτικών του παρόχου ταυτότητας, μπορεί να πλαστογραφήσει πρόσβαση σε κάθε συνδεδεμένη υπηρεσία —όπως συνέβη με την τεχνική \emph{Golden SAML} που χρησιμοποιήθηκε κατά την εκστρατεία SolarWinds. Άλλες συνήθεις αδυναμίες είναι οι εφαρμογές που δεν επικυρώνουν την υπογραφή, τον αποδέκτη ή τον χρόνο λήξης των διακριτικών· τα κλεμμένα διακριτικά συνεδρίας που επαναχρησιμοποιούνται από άλλη συσκευή· και το \emph{ψάρεμα συγκατάθεσης} (consent phishing), κατά το οποίο οι χρήστες εξαπατώνται ώστε να παραχωρήσουν ευρεία δικαιώματα OAuth σε κακόβουλη εφαρμογή. Οι άμυνες περιλαμβάνουν την προστασία των κλειδιών υπογραφής σε HSM, την αυστηρή επικύρωση των διακριτικών, τη μικρή διάρκεια ισχύος διακριτικών δεσμευμένων σε συσκευή και την έγκριση από διαχειριστή για συγκαταθέσεις OAuth υψηλού κινδύνου.

\section*{Βασικοί όροι}
\begin{description}[style=nextline,leftmargin=1.2em]
  \item[Ταυτοποίηση έναντι εξουσιοδότησης] Επαλήθευση του ποιος είναι ένα υποκείμενο έναντι απόφασης για το τι επιτρέπεται να κάνει.
  \item[FIDO2 / passkey] Ταυτοποίηση δημόσιου κλειδιού δεσμευμένη στην προέλευση, ανθεκτική στο ψάρεμα και χωρίς κοινό μυστικό.
  \item[ABAC] Έλεγχος πρόσβασης με βάση χαρακτηριστικά του υποκειμένου, του πόρου, της ενέργειας και του περιβάλλοντος.
  \item[TGT] Εισιτήριο χορήγησης εισιτηρίων του Kerberos, για τη λήψη εισιτηρίων υπηρεσίας χωρίς εκ νέου εισαγωγή διαπιστευτηρίων.
  \item[Kerberoasting] Αίτηση εισιτηρίων υπηρεσίας και «σπάσιμο» των κωδικών λογαριασμών υπηρεσιών εκτός σύνδεσης.
  \item[OIDC] OpenID Connect: επίπεδο ταυτότητας πάνω από το OAuth 2.0 που εκδίδει υπογεγραμμένα διακριτικά ταυτότητας.
\end{description}

\section*{Σύνοψη κεφαλαίου}
\begin{itemize}
  \item Η ταυτότητα είναι η σύγχρονη περίμετρος· η IAM διακρίνει την αναγνώριση, την ταυτοποίηση, την εξουσιοδότηση και την καταγραφή και διαχειρίζεται ολόκληρο τον κύκλο ζωής της ταυτότητας.
  \item Το MFA συνδυάζει παράγοντες από διαφορετικές κατηγορίες· το FIDO2 αντιστέκεται στο ψάρεμα επειδή οι υπογραφές δεσμεύονται στην προέλευση του ιστοτόπου.
  \item Το RBAC προσφέρει απλότητα, το ABAC ακρίβεια με βάση το πλαίσιο και ο έλεγχος βάσει κανόνων επιβάλλει γενικούς περιορισμούς.
  \item Το Active Directory και το Kerberos είναι στόχοι υψηλής αξίας· η κλιμακωτή διαχείριση και η σωστή διαχείριση των λογαριασμών υπηρεσιών είναι απαραίτητες.
  \item Τα πρωτόκολλα ομοσπονδίας επιτρέπουν την ενιαία σύνδεση, αλλά συγκεντρώνουν τον κίνδυνο στα κλειδιά υπογραφής και στην επικύρωση των διακριτικών.
\end{itemize}

\section*{Ερωτήσεις ανασκόπησης}
\begin{enumerate}
  \item Γιατί ένας κωδικός μίας χρήσης μέσω SMS είναι πιο αδύναμος από ένα κλειδί ασφαλείας FIDO2 απέναντι σε έναν διακομιστή ψαρέματος πραγματικού χρόνου;
  \item Σχεδιάστε μια πολιτική πρόσβασης για ένα σύστημα ιατρικών φακέλων με RBAC και, στη συνέχεια, εξειδικεύστε την με χαρακτηριστικά ABAC.
  \item Εξηγήστε γιατί το Kerberoasting είναι εφικτό για οποιονδήποτε ταυτοποιημένο χρήστη του τομέα και πώς το περιορίζουν οι λογαριασμοί gMSA.
  \item Γιατί το OAuth 2.0 από μόνο του δεν επαρκεί για την ταυτοποίηση χρηστών και τι προσθέτει το OIDC;
\end{enumerate}
